\storyheader{Khoảng lặng sau giờ tan trường
}{Ngày 09}{

Tôi ngồi trong thư viện lật từng trang sách trước mắt, chăm chú nhìn vào
từng chữ trong đó. Không gian xung quanh yên tĩnh và bình lặng tựa như
chỉ còn một mình tôi ở trên thế giới này vậy.

Tôi không thích sự ồn ào, cũng không thích sự đổi mới. Đối với tôi, chỉ
cần những giây phút yên tĩnh đọc quyển sách mình thích, bên trong thư
viện vắng tanh vì đã sau giờ tan trường, là đã trọn vẹn lắm rồi.

Cứ như vậy, ngày qua ngày tôi đều đến đây sau giờ tan trường, lựa một
quyển sách nào đó trông có vẻ thú vị và bắt đầu ngồi đọc hàng giờ liền,
ở nơi không có một bóng người nào có thể tiến đến làm phiền, không ai
gây ồn ào. Tôi cứ như vậy chìm đắm hoàn toàn vào thế giới riêng của bản
thân, đến khi sắc cam bắt đầu phủ xuống bầu trời ở bên ngoài thì tôi mới
đóng cuốn sách lại và đứng lên đi về, khóa cửa lại theo lời nhắc của cô
bạn trực thư viện.

Cuộc sống học đường của tôi cứ thế tiếp diễn như vậy. Không đổi mới,
không có gì bất ngờ, không bạn bè, không vui chơi. Mỗi ngày trôi qua tôi
đều vào thư viện và ngồi ở đó, tận hưởng từng trang chữ trong cuốn sách
với một không gian thanh bình, yên tĩnh và không ai làm phiền.

Tôi không phải một người ngại giao tiếp, lại càng không phải một kiểu
người cô độc, thế nhưng cũng không phải kiểu người sẽ có thể chủ động
bắt chuyện với ai đó. Tôi chỉ đơn giản là muốn được ở trong một thế giới
nơi tôi có thể đắm chìm vào sự yên bình mà mình thích.

Trong thế giới ấy, một mình.

Nó cũng không quá tệ.

Chỉ là...

Có hơi chút cô đơn mà thôi.

*

Như bao ngày khác, sau khi tiếng chuông tan học vang lên, không khí
trong lớp nhộn nhịp, những tiếng rủ nhau đi chơi của các nhóm bạn thân
cứ liên tục vang lên đầy háo hức, cũng có những người nằm dài ra bàn sau
một ngày học tập mệt mỏi, và cả tiếng trò chuyện vui vẻ của những người
thân thiết với nhau.

Tôi xách cặp của mình lên và đi ra khỏi lớp, không hớt hải, không hào
hứng. Thư viện không thể có chân mà chạy, và nó cũng sẽ không đóng cửa
đột xuất nên tôi cứ thể mà ung dung bước đi, vượt qua những học sinh
đang vui vẻ trò chuyện với nhau trên đường về.

Cứ như vậy, tôi đã đến thư viện của trường. Nó cách khá xa so với lớp
của tôi khi mà phải leo tận hai tầng mới đến nơi, tuy nhiên đó cũng
không phải vấn đề quá lớn, dù sao đi bộ một chút cũng tốt. Tôi kéo cánh
cửa thư viện ra, đập vào mắt là một khung cảnh quen thuộc mà tôi đã nhìn
đi nhìn lại không biết bao nhiêu lần.

Gật đầu chào cô bạn trực thư viện hôm nay, tôi tiến đến những giá sách
nằm gần đó. Tôi luôn đến thư viện mà không hề có một dự định nào trước
đó cả, chỉ đơn giản là đến kệ sách, lựa đại một quyển nào đó mà tôi thấy
trông có vẻ thú vị và ngồi đọc suốt hàng giờ đồng hồ.

Hôm nay cũng không phải ngoại lệ, tôi lang thang giữa những kệ sách lướt
qua những cuốn không hấp dẫn mấy đối với tôi. Mà thật ra cũng do tôi đã
đọc gần như toàn bộ sách ở cái thư viện này rồi nên thật sự việc tìm ra
một tác phẩm hợp ý khá là khó.

Tôi cứ thế đi lang thang trong hàng hà sa số những cuốn sách xung quanh.
Và sau đó lọt vào cặp mắt khó tính của tôi là một quyển sách với một cái
bìa trông vô cùng là đơn giản, không có họa tiết gì quá cầu kì nằm trơ
trọi giữa những cái bìa bắt mắt khác.

``Cuốn này hình như chưa đọc.'' Tôi ngắm nghía một hồi rồi quyết định
chọn nó cho buổi thư giãn ngày hôm nay. ``Lấy cái này vậy.''

Tôi đem quyển sách ra chỗ bàn trước cửa sổ. Tôi khá thích vị trí này khi
mà ánh nắng chiều tà rọi xuống từng mặt chữ khiến cảm giác đọc sách của
tôi trở nên khá dễ chịu. Bên cạnh đó thì nó cũng khá ấm áp nữa.

Tuy nhiên, cái nơi đáng lẽ ra là chỉ có mình tôi ngồi vào giờ này lại
đang có một người khác, một cô bạn cùng lớp mà tôi đã khá quen mặt tên
là Haruna.

Haruna là một người trái ngược với tôi, cậu ấy có rất nhiều bạn bè trong
lớp, điểm cao chót vót, ngoại hình đại loại là xinh đẹp. Haruna giống
như kiểu một hình mẫu lý tưởng của một học sinh con nhà người ta vậy,
điềm đạm, không quá thừa năng lượng, biết quan tâm bạn bè và học lực
xuất sắc. Nếu có điểm nào tương đồng giữa cả hai chúng tôi thì chỉ có
việc học giỏi ngang nhau mà thôi, còn lại thì cả hai cứ như hai đường
thẳng song song với nhau vậy.

Với một người nhiều bạn bè như Haruna mà sau giờ tan trường vẫn đến thư
viện thì là một việc khá là kì lạ đối với tôi. Tuy nhiên thì tôi cũng
không quá bất ngờ, thi thoảng ai đó cũng muốn đổi gió mà. Tôi định lảng
đi chỗ khác ngồi để không làm phiền đến Haruna thì bất ngờ nghe giọng
của cậu ấy vang lên.

``Cậu cứ ngồi ở đây đi, Honoka.'' Haruna nói, ánh mắt ngẩng lên nhìn vào
ghế đối diện. ``Tớ không phiền đâu.''

``Được sao?'' Tôi cất tiếng hỏi lại cậu ấy.

``Ừm.'' Haruna gật đầu xác nhận.

Tôi cũng không có vấn đề gì khi ngồi cùng cậu ấy, với lại tôi thật sự đã
quen với cảm giác đọc sách ở đây rồi nên nếu sang những bàn khác thì tôi
sẽ khá dễ mất tập trung. Nên là tôi sang ghế đối diện, ngồi xuống và lật
sách ra.

Dù nói vậy, thế nhưng tôi vẫn chả thể tập trung vào mấy con chữ bên
trong cuốn sách mà lén nhìn Haruna. Bây giờ tôi mới có dịp nhìn kĩ cậu
ấy ở khoảng cách gần như thế này, mũi cao, đôi lông mi dài, thân hình
cao ráo và gọn gàng, ánh mắt thì như muốn hút hồn người đối diện, tuy
nhiên khuôn mặt của cậu ấy bây giờ lại mang một nét nào đó trầm lắng hơn
chứ không phải sự dịu dàng mọi ngày.

Tuy vậy, tôi cũng không lén nhìn cậu ấy nữa mà quay lại với cuốn sách
trên tay mình. Thư viện đối với tôi là một nơi để bỏ hết những lo lắng
thường ngày và đắm chìm vào các trang sách, nếu Haruna đã đến đây thì
hẳn cũng không muốn bị làm phiền bởi tôi, và tôi cũng không muốn bị làm
phiền khi đang đọc sách.

Cứ như vậy, thời gian trôi đi từng giờ, tôi lẫn Haruna dán mắt vào những
trang sách được rọi bởi ánh sáng dịu nhẹ của buổi chiều, không khí yên
tĩnh phủ lên thư viện khiến tôi có thể cảm nhận rõ được sự hiện diện của
cậu ấy ở đây.

Cảm giác hôm nay khá khác với mọi khi. Nói sao nhỉ, tôi cảm thấy như
cuối cùng cũng có thể chia sẻ cái thế giới này của mình cho một người
khác, và nó khá tuyệt, tôi không biết liệu Haruna có cảm thấy như tôi
hay không, nhưng với tôi thì tôi điều này là một cảm giác khá lạ so với
ngồi đọc một mình.

*

``Đến lúc về rồi nhỉ?'' Tôi lên tiếng sau khi đã đọc xong cuốn sách dày
cộm trên tay, ngoài trời thì hoàng hôn cũng đang dần hiện lên và bầu
trời cũng đã chuyển sang sắc cam.

Cả hai đứa chúng tôi cùng đi dẹp cuốn sách của bản thân, lấy cặp và đi
ra ngoài. Tôi cũng không quên khóa cửa phòng thư viện lại, sau đó cả hai
đứa bọn tôi mang trả chìa khóa cho phòng giáo viên và cùng nhau đi về.

``Ngày nào cũng thấy cậu đến thư viện nhỉ?'' Haruna bất ngờ lên tiếng
bắt chuyện với tôi khi cả hai đang đi cùng nhau. ``Cậu thích đọc sách
lắm à?''

Cậu ấy biết sao? Bất ngờ thật, tôi là một kiểu người khá mờ nhạt trong
lớp học nên cũng không nghĩ là sẽ có thể được cô bạn siêu nổi tiếng này
chú ý đến và biết được thói quen của tôi

``Thật ra ban đầu tớ qua đó với một tâm thế là tìm một nơi yên tĩnh để
mà thư giãn thôi.'' Tôi trả lời câu hỏi của Haruna. ``Dần dần thì tớ
cũng trở nên thích đọc sách, và cũng thích luôn không khí thư giãn bên
trong thư viện.''

``Vậy sao.''

``Thế còn cậu?'' Tôi cất tiếng hỏi lại câu ban nãy của Haruna. ``Khá bất
ngờ khi cậu cũng thích đọc sách đấy, tại không mấy khi thấy cậu vào thư
viện.''

``Thích thì cũng không hẳn là thích, tớ chỉ muốn đến một nơi nào đó yên
tĩnh và thư giãn để không phải trò chuyện với ai mà thôi.'' Haruna
nghiêng đầu một chút như để suy nghĩ câu trả lời. ``Và tớ chợt nhớ đến
cậu nên quyết định vào thư viện.''

Là do tôi ấy à? Hóa ra tôi được để ý nhiều hơn tôi nghĩ ấy nhỉ?

``Nói sao nhỉ, lúc ngồi trong đó với cậu tớ cảm thấy khá thoải mái.''
Haruna nó tiếp với giọng đều đều. ``Rất yên tĩnh, nhưng không khiến tớ
cảm thấy cô đơn.''

``A, tớ cũng giống vậy.'' Tôi giơ tay lên bày tỏ ý nghĩ của mình.

Có lẽ, cả hai đứa tôi có nhiều điểm chung hơn là học lực. Đều thích cảm
giác yên tĩnh trong thư viện, đều cảm thấy khá thoải mái về đối phương
dù chưa nói chuyện quá nhiều với nhau. Hi vọng trong tương lai thì tôi
và cậu ấy sẽ thân thiết với nhau hơn.

``Ngày mai, cậu vẫn sẽ đến thư viện nhỉ?'' Tôi lên tiếng hỏi Haruna
trước khi cả hai chia tay nhau ở một ngả rẽ.

``Tất nhiên rồi, gặp lại sau giờ tan trường ngày mai nhé.''

Gặp lại sau giờ tan trường ngày mai ấy à? Tôi chưa bao giờ nghe câu tạm
biệt nào như thế này. Có lẽ trong giờ học vẫn sẽ khá khó cho bọn tôi nói
chuyện với nhau, nhưng với tôi vậy là đủ rồi.

``Vậy thì, hẹn gặp lại nhé, trong thư viện sau giờ tan trường.''
}